\documentclass[a4paper, 12pt]{article}
\usepackage[english, russian]{babel}
\usepackage[T2A]{fontenc}
\usepackage[utf8]{inputenc}
\usepackage{indentfirst}
\usepackage{geometry}
\usepackage{titleps}
\usepackage{mathtext}
\usepackage{amsmath}
\usepackage{amssymb}
\usepackage{enumitem}
\usepackage{tikz}
\setlength{\tabcolsep}{4pt}
\newpagestyle{main}{
\setheadrule{0pt}
\sethead{}{}{}
\setfootrule{0pt}
\setfoot{}{}{}
}
\pagestyle{main}

\usepackage{amsthm}
\geometry{top=2cm}

% Диаграммы Эйлера--Венна: #1 --- команды заливки, #2 --- подпись
\newcommand{\venn}[2]{%
\begin{tikzpicture}[scale=0.55]
  \begin{scope}
    #1
  \end{scope}
  \draw (0,0) circle (1) (1.2,0) circle (1);
  \node at (-0.45,0) {$A$};
  \node at (1.65,0) {$B$};
  \node at (0.6,-1.5) {#2};
\end{tikzpicture}}

\begin{document}
\begin{center}
{\large КОНСПЕКТ ЛЕКЦИЙ}

ПО ДИСЦИПЛИНЕ <<МАТЕМАТИЧЕСКИЙ АНАЛИЗ>>

\vspace{5pt}
\textbf{Элементы теории множеств и логики}
\end{center}
\vspace{10pt}

\section*{Множества}

\emph{Множество} --- произвольная совокупность определённых предметов нашей
интуиции или интеллекта, которые можно отличить один от другого и которые
представляются как единое целое.
\hfill --- Георг Кантор

Под \underline{множеством} понимают любую определённую совокупность объектов,
различимых между собой. Объекты, из которых составлено множество, ---
\underline{элементы множества}.

$A, B, X$ --- множества; \quad $a, b, x, y$ --- элементы множества.
\[
  \boxed{x \in X, \quad y \notin X}
\]

Когда строится теория множеств, используется интуитивный принцип объёмности
или аксиома экстенсиональности.

Два множества равны, если каждый элемент множества равен каждому другому
элементу множества.

Множество называется \textbf{конечным}, если оно состоит из конечного числа
элементов.

Множество называется \textbf{бесконечным}, если оно содержит бесконечное число
элементов.

Если множество не содержит ни одного элемента, то оно называется
\textbf{пустым} ($\varnothing$).

В конкретных задачах элементы всех множеств берутся из некоторого одного
достаточно широкого множества $U$, которое называется
\underline{универсальным} \underline{множеством} (универсумом).

\subsection*{Способы задания множества}

\begin{enumerate}[noitemsep]
  \item перечислением элементов;
  \item рекурсивно (перечисляющей процедурой);
  \item характеристическим свойством.
\end{enumerate}

\textbf{Пример 1 (рекурсивно).}
\begin{enumerate}[noitemsep]
  \item $3 \in X$;
  \item если $x \in X$, то $\dfrac{1}{x} \in X$ и $(1 - x) \in X$;
  \item других элементов в этом множестве нет.
\end{enumerate}
\[
  X = \left\{ 3;\ \tfrac{1}{3};\ -2;\ -\tfrac{1}{2};\ \tfrac{2}{3};\ \tfrac{3}{2} \right\}
\]

\textbf{Пример 2 (характеристическим свойством).} % в оригинале подписан как «Пример 3»
\[
  A = \{\, x : x \in \mathbb{Z},\ x^{2} \le 4 \,\} = \{-2;\ -1;\ 0;\ 1;\ 2\}
\]
\[
  B = \{\, x : x \in \mathbb{N},\ x \,\vdots\, 2 \,\} = \{2;\ 4;\ 6;\ \dots\}
\]

Чтобы избежать парадоксов, задание множеств ограничивают специальной
аксиоматикой. В частности, используют \underline{аксиому регулярности}:
\[
  X \in X \ \text{--- недопустимо!} \qquad \fbox{Парадокс брадобрея!!!}
\]

\section*{Операции над множествами}

\textbf{Соглашения:}
\begin{itemize}[noitemsep]
  \item $\Longleftrightarrow$ --- <<тогда и только тогда>>;
  \item $\exists x$ (квантор существования) --- существует $x$ такой, что\dots;
  \item $\forall x$ (квантор всеобщности) --- для любого (всякого) $x$;
  \item $\Longrightarrow$ --- следует.
\end{itemize}

Множество $A$ называется \textbf{подмножеством} множества $B$, если
$\forall x \in A\ (x \in B)$. При этом говорят, что $A$ содержится в $B$
или $B$ включает $A$:
\[
  \boxed{A \subseteq B}
\]

Если $A \subset B$ и $B \subset A$, то множества называются \textbf{равными}
($A = B$).

Множества удобно изображать с помощью кругов Эйлера (диаграмма Венна):
\begin{center}
\begin{tikzpicture}[scale=0.55]
  \draw (0,0) circle (1.6);
  \draw (-0.4,-0.2) circle (0.8);
  \node at (-0.4,-0.2) {$A$};
  \node at (0.9,0.9) {$B$};
  \node at (0,-2.1) {$A \subseteq B$};
\end{tikzpicture}
\end{center}

\textbf{Объединение множеств} $A \cup B$:
\[
  A \cup B = \{\, x : x \in A \vee x \in B \,\}
\]
\begin{center}
\venn{\fill[gray!40] (0,0) circle (1) (1.2,0) circle (1);}{$A \cup B$}
\end{center}

\textbf{Пересечение множеств} $A \cap B$ (обозначают также $AB$):
\[
  A \cap B = \{\, x : x \in A \wedge x \in B \,\}
\]
\begin{center}
\venn{\clip (0,0) circle (1); \fill[gray!40] (1.2,0) circle (1);}{$A \cap B$}
\end{center}

\textbf{Разность множеств} $A \setminus B$:
\begin{center}
\venn{\clip (0,0) circle (1); \fill[gray!40, even odd rule] (0,0) circle (1) (1.2,0) circle (1);}{$A \setminus B$}
\qquad
\venn{\clip (1.2,0) circle (1); \fill[gray!40, even odd rule] (0,0) circle (1) (1.2,0) circle (1);}{$B \setminus A$}
\end{center}

\textbf{Симметричная разность множеств}:
\[
  A \,\triangle\, B = (A \setminus B) \cup (B \setminus A)
\]
\begin{center}
\venn{\fill[gray!40, even odd rule] (0,0) circle (1) (1.2,0) circle (1);}{$A \,\triangle\, B$}
\end{center}

\textbf{Дополнение множества} относительно универсального множества:
\[
  \overline{A} = U \setminus A = CA
\]
\begin{center}
\begin{tikzpicture}[scale=0.55]
  \fill[gray!40] (-2,-1.4) rectangle (2,1.4);
  \fill[white] (0,0) circle (0.8);
  \draw (-2,-1.4) rectangle (2,1.4);
  \draw (0,0) circle (0.8);
  \node at (0,0) {$A$};
  \node at (1.6,1.05) {$U$};
  \node at (0,-1.9) {$\overline{A}$};
\end{tikzpicture}
\end{center}

\subsection*{Свойства операций}

\begin{itemize}[noitemsep]
  \item \textbf{коммутативность}:
    \[ A \cup B = B \cup A, \qquad A \cap B = B \cap A; \]
  \item \textbf{ассоциативность}:
    \[ (A \cup B) \cup C = A \cup (B \cup C), \qquad
       (A \cap B) \cap C = A \cap (B \cap C); \]
  \item \textbf{дистрибутивность}:
    \[ (A \cup B) \cap C = (A \cap C) \cup (B \cap C), \qquad
       (A \cap B) \cup C = (A \cup C) \cap (B \cup C); \]
  \item \textbf{идемпотентность}:
    \[ A \cup A = A, \qquad A \cap A = A; \]
  \item \textbf{законы де Моргана} (правила двойственности):
    \[ \overline{A \cup B} = \overline{A} \cap \overline{B}, \qquad
       \overline{A \cap B} = \overline{A} \cup \overline{B}; \]
  \item \textbf{поглощение}:
    \[ A \cup (A \cap B) = A, \qquad A \cap (A \cup B) = A; \]
  \item свойства $\overline{A}$, $\varnothing$ и $U$: % заголовок пункта в оригинале зачёркнут
    \[
      \begin{aligned}
        A \cup \overline{A} &= U, & A \cup U &= U, \\
        A \cap \overline{A} &= \varnothing, & A \cap U &= A, \\
        A \cup \varnothing &= A, & U \cup \varnothing &= U, \\
        A \cap \varnothing &= \varnothing, & U \cap \varnothing &= \varnothing.
      \end{aligned}
    \]
\end{itemize}

\section*{Алгебра высказываний}

\textbf{Высказывание} --- такое предложение, которое либо истинно, либо ложно.
Высказывание не может быть одновременно и истинным, и ложным.

Высказывания обозначают заглавными буквами.

Сопоставив истинному высказыванию символ $1$, а ложному --- $0$, введём
функцию $\lambda$, заданную на совокупности всех высказываний и отвечающую
следующему требованию:
\[
  \lambda(P) =
  \begin{cases}
    1, & \text{если } P \text{ --- истинно};\\
    0, & \text{если } P \text{ --- ложно}.
  \end{cases}
\]
Функцию $\lambda$ называют \textbf{функцией истинности}, а значение
$\lambda(P)$ --- \textbf{логическим значением} высказывания $P$.

\begin{itemize}
  \item \textbf{Отрицанием} высказывания $P$ называется новое высказывание,
    обозначаемое $\neg P$, которое истинно, если высказывание $P$ ложно,
    и ложно, если высказывание $P$ истинно.
    \[ \boxed{\neg P,\quad \overline{P},\quad !P} \]

  \item \textbf{Конъюнкцией} двух высказываний $P$ и $Q$ называется новое
    высказывание, обозначаемое $P \wedge Q$ или $P \,\&\, Q$, которое истинно
    лишь в единственном случае, когда истинны оба высказывания $P$ и $Q$,
    и ложно во всех остальных случаях (логическое~$\times$).
    \begin{center}
    \begin{tabular}{c|c|c}
      $\lambda(P)$ & $\lambda(Q)$ & $\lambda(P \wedge Q)$ \\ \hline
      0 & 0 & 0 \\
      0 & 1 & 0 \\
      1 & 0 & 0 \\
      1 & 1 & 1 \\
    \end{tabular}
    \end{center}

  \item \textbf{Дизъюнкцией} двух высказываний $P$ и $Q$ называется новое
    высказывание, обозначаемое $P \vee Q$, которое истинно тогда и только
    тогда, когда истинно хотя бы одно из высказываний $P$ или $Q$, и ложно
    в единственном случае, когда оба высказывания $P$ и $Q$ ложны
    (логическое~$+$).
    \begin{center}
    \begin{tabular}{c|c|c}
      $\lambda(P)$ & $\lambda(Q)$ & $\lambda(P \vee Q)$ \\ \hline
      0 & 0 & 0 \\
      0 & 1 & 1 \\
      1 & 0 & 1 \\
      1 & 1 & 1 \\
    \end{tabular}
    \end{center}

  \item \textbf{Импликацией} двух высказываний $P$ и $Q$ называется новое
    высказывание, обозначаемое $P \Rightarrow Q$ (<<если $P$, то $Q$>>,
    <<из $P$ следует $Q$>>), которое ложно в единственном случае, когда
    высказывание $P$ истинно, а $Q$ ложно, а во всех остальных случаях ---
    истинно.
    \begin{center}
    \begin{tabular}{c|c|c}
      $\lambda(P)$ & $\lambda(Q)$ & $\lambda(P \Rightarrow Q)$ \\ \hline
      0 & 0 & 1 \\
      0 & 1 & 1 \\
      1 & 0 & 0 \\
      1 & 1 & 1 \\
    \end{tabular}
    \end{center}

  \item \textbf{Эквиваленцией} двух высказываний $P$ и $Q$ называется новое
    высказывание $P \Leftrightarrow Q$ (<<тогда и только тогда>>), которое
    истинно в том и только в том случае, когда оба высказывания $P$ и $Q$
    одновременно истинны или оба ложны, а во всех остальных случаях --- ложно.
    \begin{center}
    \begin{tabular}{c|c|c}
      $\lambda(P)$ & $\lambda(Q)$ & $\lambda(P \Leftrightarrow Q)$ \\ \hline
      0 & 0 & 1 \\
      0 & 1 & 0 \\
      1 & 0 & 0 \\
      1 & 1 & 1 \\
    \end{tabular}
    \end{center}
\end{itemize}

\textbf{!} Операции над высказываниями аналогичны операциям над множествами.
Свойства логических операций доказывают с помощью таблиц истинности.

\end{document}
